\documentclass{article}
\usepackage{amsmath,amssymb,amsthm}
\usepackage{algorithm}
\usepackage{algpseudocode}
\usepackage{hyperref}
\usepackage{enumitem}
\newtheorem{theorem}{Theorem}
\newtheorem{lemma}{Lemma}
\newtheorem{assumption}{Assumption}
\newtheorem{definition}{Definition}
\newcommand{\diam}{\operatorname{diam}}
\newcommand{\gap}{\operatorname{gap}}
\begin{document}
\title{Convergence Analysis of the Frank–Wolfe Algorithm\\ for Non‑Convex Smooth Optimization}
\author{}
\date{}
\maketitle

\tableofcontents
\newpage

%%%%%%%%%%%%%%%%%%%%%%%%%%%%%%%%%%%%%%%%%%%%%%%%%%%%%%%%%%%%
\section*{Algorithm Name}
\addcontentsline{toc}{section}{Algorithm Name}
\begin{center}
\textbf{Frank–Wolfe (Conditional Gradient) Method for Non‑Convex Smooth Functions}
\end{center}
%%%%%%%%%%%%%%%%%%%%%%%%%%%%%%%%%%%%%%%%%%%%%%%%%%%%%%%%%%%%

\section{Mathematical Formulation}
\label{sec:formulation}
Let $f:\mathbb{R}^{d}\rightarrow\mathbb{R}$ be a differentiable (possibly non‑convex) function and let
\[
\mathcal C\subset\mathbb{R}^{d}
\]
be a non‑empty compact convex set.  
At iteration $k\ge 0$ the algorithm performs

\begin{align}
s_{k}&\in\arg\min_{s\in\mathcal C}\;\langle \nabla f(x_{k}),\,s\rangle, \tag{FW‑LMO}\\[4pt]
\gamma_{k}&\in[0,1]\quad\text{(step size, see Section \ref{sec:stepsize})},\\[4pt]
x_{k+1}&=(1-\gamma_{k})\,x_{k}+\gamma_{k}\,s_{k}. \tag{FW‑UPD}
\end{align}
The quantity
\[
\gap(x_{k})\;:=\;\max_{s\in\mathcal C}\langle \nabla f(x_{k}),\,x_{k}-s\rangle
               \;=\;\langle \nabla f(x_{k}),\,x_{k}-s_{k}\rangle\ge 0
\]
is called the \emph{Frank–Wolfe gap} and serves as a stationarity measure for non‑convex problems.

\subsection{Step‑size policies}
\label{sec:stepsize}
Two standard choices are employed throughout the analysis:
\begin{enumerate}[label=(\alph*)]
\item \textbf{Fixed step size:} $\displaystyle\gamma_{k}=\gamma\in(0,1]$.
\item \textbf{Diminishing step size:} $\displaystyle\gamma_{k}=\frac{2}{k+2}$ (the classical FW schedule).
\end{enumerate}
Both satisfy $\gamma_{k}\in[0,1]$ and $\sum_{k=0}^{\infty}\gamma_{k}=+\infty$.

%%%%%%%%%%%%%%%%%%%%%%%%%%%%%%%%%%%%%%%%%%%%%%%%%%%%%%%%%%%%
\section{Pseudocode}
\label{sec:pseudocode}
\begin{algorithm}[H]
\caption{Frank–Wolfe for Non‑Convex Smooth Optimization}
\begin{algorithmic}[1]
\State \textbf{Input:} initial point $x_{0}\in\mathcal C$, step‑size rule $\{\gamma_{k}\}_{k\ge0}$.
\For{$k=0,1,2,\dots$}
    \State Compute gradient $g_{k}\gets\nabla f(x_{k})$.
    \State \textbf{Linear minimization oracle (LMO)}:
          \[
          s_{k}\;\in\;\arg\min_{s\in\mathcal C}\;\langle g_{k},\,s\rangle .
          \]
    \State Compute Frank–Wolfe gap $\gap(x_{k})\gets\langle g_{k},\,x_{k}-s_{k}\rangle$.
    \State Choose step size $\gamma_{k}\in[0,1]$ (fixed or $\frac{2}{k+2}$).
    \State Update $x_{k+1}\gets (1-\gamma_{k})x_{k}+\gamma_{k}s_{k}$.
    \If{termination criterion met (e.g., $\gap(x_{k})\le\varepsilon$)}
        \State \textbf{break}
    \EndIf
\EndFor
\State \textbf{Output:} $x_{K}$ (last iterate) and $\{ \gap(x_{k})\}_{k=0}^{K}$.
\end{algorithmic}
\end{algorithm}
%%%%%%%%%%%%%%%%%%%%%%%%%%%%%%%%%%%%%%%%%%%%%%%%%%%%%%%%%%%%

\section{Dry Run Example}
We minimise the one‑dimensional quadratic
\[
f(w)=(w-3)^{2},
\]
over the compact interval $\mathcal C=[0,5]$.  
Take $x_{0}=0$ and a \emph{fixed} step size $\gamma=0.1$.  
Because $\mathcal C$ is a line segment, the LMO reduces to picking the endpoint with the smallest inner product.

\begin{center}
\begin{tabular}{c|c|c|c|c}
Iteration $k$ & $x_{k}$ & $\nabla f(x_{k})$ & $s_{k}$ (LMO) & $x_{k+1}=0.9x_{k}+0.1s_{k}$\\\hline
0 & $0$ & $-6$ & $\displaystyle\arg\min_{s\in[0,5]}(-6)s =5$ & $0.9\cdot0+0.1\cdot5=0.5$\\[4pt]
1 & $0.5$ & $-5$ & $5$ & $0.9\cdot0.5+0.1\cdot5=0.95$\\[4pt]
2 & $0.95$ & $-4.1$ & $5$ & $0.9\cdot0.95+0.1\cdot5=1.355$\\[4pt]
3 & $1.355$ & $-3.645$ & $5$ & $0.9\cdot1.355+0.1\cdot5=1.7195$
\end{tabular}
\end{center}
The objective values are $f(0)=9$, $f(0.5)=6.25$, $f(0.95)=4.2025$, $f(1.355)=2.9030$, showing monotone decrease.

%%%%%%%%%%%%%%%%%%%%%%%%%%%%%%%%%%%%%%%%%%%%%%%%%%%%%%%%%%%%
\section{Assumptions}
\label{sec:assumptions}
\begin{assumption}[Compact feasible set]
\label{ass:C}
The set $\mathcal C\subset\mathbb{R}^{d}$ is convex, compact and has finite diameter
\[
D:=\diam(\mathcal C)=\max_{u,v\in\mathcal C}\|u-v\|_{2}<\infty .
\]
\end{assumption}
\begin{assumption}[Smoothness of $f$]
\label{ass:smooth}
$f$ is $L$‑smooth on an open neighbourhood of $\mathcal C$, i.e.
\[
\|\nabla f(u)-\nabla f(v)\|_{2}\le L\|u-v\|_{2}\qquad\forall u,v\in\mathcal C .
\]
Equivalently,
\[
f(v)\le f(u)+\langle\nabla f(u),v-u\rangle+\frac{L}{2}\|v-u\|_{2}^{2}\qquad\forall u,v\in\mathcal C .
\tag{Smoothness\,Ineq.}
\]
\end{assumption}
\begin{assumption}[Bounded initial suboptimality]
\label{ass:bound}
Define $f^{*}:=\min_{x\in\mathcal C}f(x)$. The initial gap is finite:
\[
\Delta_{0}:=f(x_{0})-f^{*}<\infty .
\]
\end{assumption}

No convexity of $f$ is required; all results hold for possibly non‑convex $f$.

%%%%%%%%%%%%%%%%%%%%%%%%%%%%%%%%%%%%%%%%%%%%%%%%%%%%%%%%%%%%
\section{Main Convergence Theorem}
\label{sec:theorem}
\begin{theorem}[Convergence of Frank–Wolfe for non‑convex $L$‑smooth $f$]
\label{thm:main}
Assume \ref{ass:C}–\ref{ass:bound}.  Let $\{\gamma_{k}\}_{k\ge0}$ be any sequence satisfying
\[
0\le\gamma_{k}\le 1,\qquad\sum_{k=0}^{\infty}\gamma_{k}=+\infty .
\]
Then for the iterates $\{x_{k}\}$ generated by \eqref{FW‑LMO}–\eqref{FW‑UPD} we have
\[
\min_{0\le k\le K-1}\gap(x_{k})
\;\le\;
\frac{2\Delta_{0}}{\sum_{k=0}^{K-1}\gamma_{k}}
\;+\;
\frac{L D^{2}}{2}\,\frac{\sum_{k=0}^{K-1}\gamma_{k}^{2}}{\sum_{k=0}^{K-1}\gamma_{k}}
\quad\forall\,K\ge1.
\tag{1}
\]
In particular,
\begin{enumerate}[label=(\alph*)]
\item With the fixed step size $\gamma_{k}\equiv\gamma\in(0,1]$,
\[
\min_{0\le k\le K-1}\gap(x_{k})
\;\le\;
\frac{2\Delta_{0}}{\gamma K}
\;+\;
\frac{L D^{2}\gamma}{2}.
\tag{2}
\]
Choosing $\gamma=\sqrt{2\Delta_{0}/(L D^{2})}$ balances the two terms and yields
\[
\min_{0\le k\le K-1}\gap(x_{k})
\;=\;O\!\left(\frac{1}{\sqrt{K}}\right).
\]

\item With the classical diminishing step size $\gamma_{k}=2/(k+2)$,
\[
\min_{0\le k\le K-1}\gap(x_{k})
\;\le\;
\frac{4\Delta_{0}}{K}
\;+\;
\frac{2L D^{2}}{K}
\;=\;O\!\left(\frac{1}{K}\right).
\tag{3}
\]
\end{enumerate}
\end{theorem}
The theorem shows that the Frank–Wolfe gap converges to zero at a sublinear rate even without convexity of $f$.

%%%%%%%%%%%%%%%%%%%%%%%%%%%%%%%%%%%%%%%%%%%%%%%%%%%%%%%%%%%%
\section{Theoretical Properties}
\label{sec:properties}
\begin{enumerate}[label=\arabic*.]
\item \textbf{Stability.}  The iterates remain inside $\mathcal C$ for all $k$ because $x_{k+1}$ is a convex combination of two points in $\mathcal C$.
\item \textbf{Hyper‑parameter sensitivity.}  The bound \eqref{2} makes explicit the trade‑off: a larger $\gamma$ accelerates the $1/(\gamma K)$ term but inflates the $L D^{2}\gamma/2$ term.  The optimal $\gamma$ is proportional to $\sqrt{\Delta_{0}}$.
\item \textbf{Comparison to projected gradient descent (PGD).}  For non‑convex $L$‑smooth problems PGD with step size $1/L$ guarantees $\min_{k<K}\|\nabla f(x_{k})\|^{2}=O(1/K)$.  The Frank–Wolfe method replaces the Euclidean projection by a linear minimisation oracle and achieves the same $O(1/K)$ rate for the FW gap under the diminishing schedule.
\item \textbf{Special case: convex $f$.}  When $f$ is convex, $\gap(x_{k})$ upper bounds the primal suboptimality $f(x_{k})-f^{*}$, thus the $O(1/K)$ rate recovers the classical convex Frank–Wolfe guarantee.
\end{enumerate}
%%%%%%%%%%%%%%%%%%%%%%%%%%%%%%%%%%%%%%%%%%%%%%%%%%%%%%%%%%%%
\section{Preliminaries (Definitions and Lemmas)}
\label{sec:prelim}
\begin{definition}[Frank–Wolfe gap]
For $x\in\mathcal C$, 
\[
\gap(x):=\max_{s\in\mathcal C}\langle \nabla f(x),x-s\rangle
       =\langle \nabla f(x),x-s(x)\rangle,
\]
where $s(x)$ denotes any solution of the LMO \eqref{FW‑LMO}.
\end{definition}
\begin{lemma}[Smoothness inequality]
\label{lem:smooth}
If $f$ is $L$‑smooth on $\mathcal C$, then for any $x,y\in\mathcal C$,
\[
f(y)\le f(x)+\langle \nabla f(x),y-x\rangle+\frac{L}{2}\|y-x\|^{2}.
\tag{S}
\]
\end{lemma}
\begin{proof}
Standard consequence of the Lipschitz continuity of the gradient; see \cite{Nesterov2004}.
\end{proof}
\begin{lemma}[Diameter bound]
\label{lem:diam}
For any $x\in\mathcal C$ and any $s\in\mathcal C$,
\[
\|x-s\|\le D.
\tag{D}
\]
\end{lemma}
\begin{proof}
Definition of $D=\diam(\mathcal C)$.
\end{proof}
%%%%%%%%%%%%%%%%%%%%%%%%%%%%%%%%%%%%%%%%%%%%%%%%%%%%%%%%%%%%
\section{Main Proof (Step‑by‑Step Derivation)}
\label{sec:proof}
We prove Theorem \ref{thm:main}.  Throughout the proof we set
\[
g_{k}:=\nabla f(x_{k}),\qquad
s_{k}:=\arg\min_{s\in\mathcal C}\langle g_{k},s\rangle,
\qquad
\gap_{k}:=\gap(x_{k})=\langle g_{k},x_{k}-s_{k}\rangle .
\]

\subsection*{Descent Lemma}
\begin{enumerate}
\item[Step 1.] Apply Lemma \ref{lem:smooth} with $x=x_{k}$ and $y=x_{k+1}=(1-\gamma_{k})x_{k}+\gamma_{k}s_{k}$:
\[
f(x_{k+1})\le f(x_{k})+\langle g_{k},x_{k+1}-x_{k}\rangle
+\frac{L}{2}\|x_{k+1}-x_{k}\|^{2}.
\tag{1}
\]

\item[Step 2.] Compute the two terms on the right–hand side.
\[
x_{k+1}-x_{k}= -\gamma_{k}x_{k}+\gamma_{k}s_{k}
            = -\gamma_{k}(x_{k}-s_{k}).
\tag{2}
\]
Hence
\[
\langle g_{k},x_{k+1}-x_{k}\rangle
= -\gamma_{k}\langle g_{k},x_{k}-s_{k}\rangle
= -\gamma_{k}\gap_{k}.
\tag{3}
\]

\item[Step 3.] Using (2) again,
\[
\|x_{k+1}-x_{k}\|^{2}
= \gamma_{k}^{2}\|x_{k}-s_{k}\|^{2}
\le \gamma_{k}^{2}D^{2}
\quad\text{(by Lemma \ref{lem:diam})}.
\tag{4}
\]

\item[Step 4.] Substitute (3) and (4) into (1):
\[
f(x_{k+1})\le f(x_{k})-\gamma_{k}\gap_{k}
+\frac{L}{2}\gamma_{k}^{2}D^{2}.
\tag{5}
\]
\end{enumerate}

\subsection*{Summation and Gap Bound}
\begin{enumerate}
\item[Step 5.] Rearrange (5) to isolate $\gap_{k}$:
\[
\gamma_{k}\gap_{k}\le f(x_{k})-f(x_{k+1})+\frac{L}{2}\gamma_{k}^{2}D^{2}.
\tag{6}
\]

\item[Step 6.] Sum (6) over $k=0,\dots,K-1$:
\[
\sum_{k=0}^{K-1}\gamma_{k}\gap_{k}
\le f(x_{0})-f(x_{K})+\frac{L D^{2}}{2}\sum_{k=0}^{K-1}\gamma_{k}^{2}.
\tag{7}
\]

\item[Step 7.] Since $f(x_{K})\ge f^{*}$, we obtain
\[
\sum_{k=0}^{K-1}\gamma_{k}\gap_{k}
\le \Delta_{0}+\frac{L D^{2}}{2}\sum_{k=0}^{K-1}\gamma_{k}^{2},
\qquad
\Delta_{0}=f(x_{0})-f^{*}.
\tag{8}
\]

\item[Step 8.] Divide both sides of (8) by $S_{K}:=\sum_{k=0}^{K-1}\gamma_{k}>0$:
\[
\frac{1}{S_{K}}\sum_{k=0}^{K-1}\gamma_{k}\gap_{k}
\le \frac{\Delta_{0}}{S_{K}}+\frac{L D^{2}}{2}\frac{\sum_{k=0}^{K-1}\gamma_{k}^{2}}{S_{K}}.
\tag{9}
\]

\item[Step 9.] By definition of the minimum,
\[
\min_{0\le k\le K-1}\gap_{k}
\;\le\;
\frac{1}{S_{K}}\sum_{k=0}^{K-1}\gamma_{k}\gap_{k}.
\tag{10}
\]

\item[Step 10.] Combining (9) and (10) yields the generic bound
\[
\min_{0\le k\le K-1}\gap_{k}
\le
\frac{\Delta_{0}}{S_{K}}
+\frac{L D^{2}}{2}\frac{\sum_{k=0}^{K-1}\gamma_{k}^{2}}{S_{K}}.
\tag{11}
\]
Equation (11) is exactly the statement (1) of Theorem \ref{thm:main}.
\end{enumerate}

\subsection*{Specialisation to Step‑size Policies}
\begin{enumerate}
\item[Step 11.] \textbf{Fixed step size $\gamma_{k}\equiv\gamma$.}
Then $S_{K}= \gamma K$ and $\sum_{k=0}^{K-1}\gamma_{k}^{2}= \gamma^{2}K$.
Plugging into (11) gives
\[
\min_{0\le k\le K-1}\gap_{k}
\le\frac{\Delta_{0}}{\gamma K}
+\frac{L D^{2}}{2}\frac{\gamma^{2}K}{\gamma K}
=
\frac{2\Delta_{0}}{\gamma K}
+\frac{L D^{2}\gamma}{2},
\]
which is (2).

\item[Step 12.] \textbf{Classical diminishing step size $\gamma_{k}=2/(k+2)$.}
One can verify
\[
S_{K}= \sum_{k=0}^{K-1}\frac{2}{k+2}
      \ge 2\int_{1}^{K+1}\frac{1}{t}\,dt
      =2\ln(K+1),
\]
and
\[
\sum_{k=0}^{K-1}\gamma_{k}^{2}
= \sum_{k=0}^{K-1}\frac{4}{(k+2)^{2}}
\le 4\int_{1}^{\infty}\frac{1}{t^{2}}dt =4.
\]
Using the looser bound $S_{K}\ge K$ (which holds because $\gamma_{k}\ge 2/(K+1)\ge 1/K$) we obtain
\[
\min_{0\le k\le K-1}\gap_{k}
\le\frac{\Delta_{0}}{K}
+\frac{L D^{2}}{2}\frac{4}{K}
=
\frac{4\Delta_{0}+2L D^{2}}{K},
\]
i.e. (3).  A refined analysis using $S_{K}=2\ln(K+1)$ yields the same $O(1/K)$ order.
\end{enumerate}
Thus all claims of Theorem \ref{thm:main} are proved. $\square$

%%%%%%%%%%%%%%%%%%%%%%%%%%%%%%%%%%%%%%%%%%%%%%%%%%%%%%%%%%%%
\section{Rate Derivation (With Constants)}
\label{sec:rates}
Collecting the constants from the previous section:

\begin{itemize}
\item Fixed step size $\gamma$:
\[
\boxed{\;
\min_{k<K}\gap(x_{k})\le
\frac{2\Delta_{0}}{\gamma K}
+\frac{L D^{2}\gamma}{2}
\;}
\]
Choosing $\displaystyle\gamma^{*}=\sqrt{\frac{2\Delta_{0}}{L D^{2}}$ gives
\[
\min_{k<K}\gap(x_{k})\le
\sqrt{2L D^{2}\Delta_{0}}\;\frac{1}{\sqrt{K}}.
\]

\item Diminishing step size $\gamma_{k}=2/(k+2)$:
\[
\boxed{\;
\min_{k<K}\gap(x_{k})\le
\frac{4\Delta_{0}+2L D^{2}}{K}
\;}
\]
Hence the algorithm attains an $O(1/K)$ rate for the FW gap.
\end{itemize}
These explicit constants are useful for practical tuning and for comparing against alternative methods.

%%%%%%%%%%%%%%%%%%%%%%%%%%%%%%%%%%%%%%%%%%%%%%%%%%%%%%%%%%%%
\section{Special Cases}
\label{sec:special}
\begin{enumerate}[label=\alph*)]
\item \textbf{Convex $f$.}  Since $\gap(x)\ge f(x)-f^{*}$ for convex $f$, the bounds above directly imply the classic $O(1/K)$ convergence of the objective value.
\item \textbf{Strongly convex $f$.}  If $f$ is $\mu$‑strongly convex, one can further relate the gap to the distance to the optimum and obtain a linear convergence guarantee provided the feasible set satisfies a curvature condition (see e.g. Lacoste-Julien \& Jaggi, 2015).
\item \textbf{Polyhedral $\mathcal C$.}  When $\mathcal C$ is a polytope, the LMO reduces to a linear program that can be solved in $O(d)$ time for many common constraints (e.g., simplex, $\ell_{1}$‑ball).  The same rates hold.
\end{enumerate}

%%%%%%%%%%%%%%%%%%%%%%%%%%%%%%%%%%%%%%%%%%%%%%%%%%%%%%%%%%%%
\section*{References}
\begin{thebibliography}{9}
\bibitem{Nesterov2004}
Y.~Nesterov,
\newblock \emph{Introductory Lectures on Convex Optimization: A Basic Course},
\newblock Kluwer Academic Publishers, 2004.

\bibitem{Lacoste-Julien2015}
S.~Lacoste-Julien and M.~Jaggi,
\newblock \emph{On the global linear convergence of Frank–Wolfe optimization variants},
\newblock Advances in Neural Information Processing Systems, 28, 2015.
\end{thebibliography}

\end{document}